\section{\ftwoz \  for UPESAT}
%\john{Can't we operate on multiple polynomial rings, namely $Z[X]$ and $F_2[X]$} \albert{Yes but we don't want the complexity of doing so}



\subsection{Universal Polynomial Equation Satisfiability (UPESAT)}

%\albertimportant{todo: add the virtualization formalism. }

%Throughout this section $R$ denotes either $\ZZ$ or a finite field $\FF$ with $q^e$ elements. 
\albertimportant{[No time for EC?] It may be good to allow $R$ to be $\QQ$, so we can inline constraints more easily (even though we can write the constraints over $\QQ$ and then turn them into $\ZZ$ constraints by multiplying by the LCM, so maybe it's not worth the added complexity}

In this section we describe a variation of the PESAT  relation (cf.\ \cref{d:pesat}), called \emph{Universal Polynomial Equation Satisfiability (UPESAT)}. It generalizes PESAT by allowing polynomial constraints over a collection of quotient rings of the polynomial ring $\ring[X]$. UPESAT is a subset of the UCS relation from \cite{zincplus} (cf.\ \cref{r:pimsat_vs_ucs}). %On the other hand, UPESAT operates exclusively on the polynomial ring $\ZZ[X]$.



 
%We define a ring homomorphism $\pi:\ZZ \to \ring$ as follows 
%\begin{align*}
%\pi: \ZZ &\to \ring\\
%x &\mapsto \begin{cases}
%x & \text{if } \ring=\ZZ,\\
%\pi_q(x) & \text{if } \ring=\FF,
%\end{cases}
%\end{align*}
%where, recall $\pi_q$ denotes reduction modulo $q$.  We use the same notation for the $0$ and $1$ element of $\ring$ as the integers $0$ and $1$, but we will still use the homomorphism $\pi$ and its inverse to map between the two \albert{comment earlier in prelims?}.



We fix a witness length $\codedim_0$ which we assume is a power of two; an input length $\inplen$; a bit size bound $B$; a degree bound $\moduledegree$; a collection of $\iota$ rings $R_1, \ldots, R_{\iota}$ of \emph{zero or prime characteristic} (our constraints will be over these rings). As all rings in this paper, we assume $R_1, \ldots, R_\iota$ are finitely generated as rings (see \cref{s:prelims_algebra}). We fix a number of rows $\numrows$; a ring $S$ (our witness will have entries in $S$); and for each $i\in [\iota]$, a   ring homomorphisms $\psi_i: S\to R_i$ and a  vector of $\numrows$ polynomials
$$\mathbf{p}^{(i)} = \left(\left.\mathbf{p}_{j}^{(i)} \ \right| \ j\in[\numrows]\right) \in \left(R_i[Y_1,\ldots, Y_{\inplen+\codedim}]\right)^{\numrows}$$ with coefficients in $\ring_i$ and on $\inplen + \codedim$ variables. %Given $a\in \ring_i$ and $i\in [\iota]$, we say that \emph{$a=0$ holds in $R_i$} if the image of $a$ under the canonical projection $\ring[X]\to R_i$ is zero. 
We denote $$\bm{\psi}=(\psi_i)_{i\in [\iota]},\qquad \pp = (\pp^{(i)})_{i\in [\iota]}, \qquad \bm{R}=(R_i)_{i\in[\iota]}.$$ 
%
The UPESAT relation with respect to these parameters is defined as
$$
%\UPESAT(\ringtwobounds{B}{\moduledegree}, \pp,\bm{R}, \cC, \pack) 
\UPESAT(\moduleonebound, \bm{\psi}, \pp,\bm{R}) 
= \left\{\left. 
\begin{aligned}
    \inp &= \xx \in \left(\moduleonebound\right)^{\inplen},\\
    \oracle &=\bot,\\
    %\oracle &= f \in \comfield^{\codelength},\\
    \wit &=\bff \in \left(\moduleonebound\right)^{\codedim_0}
\end{aligned} 
\ \right| \ 
\begin{aligned}
    %& f = \Enc_{\cC}\circ \pack \circ \picomfield(\bitstext(\bff)),\\
    &\mathbf{p}_{j}^{(i)}(\psi_i(\xx), \psi_i(\bff)) = 0 \ \text{ in } R_i \ \text{ for all } j\in [\numrows], \ i \in [\iota]
\end{aligned}
\right\}
$$
(See \cref{s:prelims_algebra} for $\moduleonebound$). 
%For instance, each $R_i$ can be $\ring[X]$ itself (taking $h_i=0$), the base ring $\ring$ (taking $h_i=X$), a cyclotomic ring such as $\ZZ[X]/(X^{64}+1)$, or, when $\ring=\FF$, any finite extension field of $\FF$ (taking $h_i$ irreducible over $\FF$). When $h_i=0$ for all $i\in [\iota]$, so that the constraints are plain polynomial equations over $\ring[X]$, we write $\UPESAT(\ringtwobounds{B}{\moduledegree}, \pp,\bot)$.
%When $\ideal_i=0$ for all $i$, then $\UPESAT(\ringtwobounds{B}{\moduledegree}, \pp, \bm{R}) = \PESAT(\)$

%The next remark explains that all our constructions for $\UPESAT$ apply, in particular, to the usual R1CS, CCS \cite{ccs}, AIR \cite{ethstark}, lookups,  permutation constraints, etc.\ over finite fields.


Following  \cref{e:pimsat_ronecs} and \cite{linear-accumulation,zinc,zincplus} one can  show that UPESAT allows to express  popular constraint systems such as R1CS, CCS, AIR, and lookup and permutation constraints, over different rings $R_i$, with witnesses and public inputs having entries in a ring $S$ that projects to the rings $R_i$.


\begin{example}[R1CS constraints over quotients of {$\ring[X]$}]\label{e:pimsat_ronecs}
Let $S,\inplen,\codedim_0,\numrows,\iota, \bm{R}, \bm{\psi}, \pp, B$ be as before,  with $R_1=\ldots = R_\iota=S$, $\psi_i=\mathsf{Id}$ for all $i$, and let $A^{(i)},B^{(i)},C^{(i)}$ be $\numrows \times (\inplen+\codedim)$ matrices with entries in $S$ (typically, these are sparse matrices with entries in a small subset of $S$), for all $i\in [\iota]$. Then the R1CS-UPESAT relation is defined as:
    $$
\ronecsideal \defeq \left\{\left. 
\begin{aligned}
    \inp &= \xx\in \left(\moduleonebound\right)^{\inplen},\\
    \oracle &= \bot,\\
    %\oracle &= f \in \comfield^{\codelength},\\
    \wit &=\bff \in \left(\moduleonebound\right)^{\codedim_0}
\end{aligned} 
\ \right| \ 
\begin{aligned}
 %& f = \Enc_{\cC}\circ \pack \circ \picomfield(\bitstext(\bff)),\\  
 &\left(\left(A^{(i)}\cdot \zz\right)\circ \left(B^{(i)}\cdot \zz\right) - \left(C^{(i)}\cdot \zz\right)\right)_{j} = 0 \ \text{ in } S \\ &\qquad\qquad \qquad \qquad \qquad \text{ for all } j\in [\numrows], \ i \in [\iota]\\
    &\ \zz=(\xx, \bff)
\end{aligned}
\right\},
$$
The relation $\ronecsideal$ is  precisely $\UPESAT(S,\bm{\psi}, \pp,\bm{R})$ with $$\mathbf{p}_{j}^{(i)}(\xx,\bff)= \left(\left(A^{(i)}\cdot \zz\right)\circ \left(B^{(i)}\cdot \zz\right) - \left(C^{(i)}\cdot \zz\right)\right)_{j}.$$
By selecting, say $S$ a finite field, one recovers the standard R1CS constraint. Setting $S=\ZZ$ one obtains an R1CS constraint over the integers.
\end{example}
%



\begin{remark}[UPESAT vs the UCS relation from Zinc+ \cite{zincplus}]\label{r:pimsat_vs_ucs}
UPESAT is a closely related to the  UCS relation from \cite{zincplus}. Fixing $S=Z[X](X^d)$ for some $d$ leads to UPESAT being a large subset of UCS. In this case, the only  ways in which UCS is more general is that: 1) it supports witnesses with entries in rings of the form  $\QQ[X]/(X^d)$; 2) it allows to have several witness vectors, with one of them typed over $S$, and others typed over the rings $R_i$ (as opposed to  having a unique witness vector typed in $S$).

Other than that, UPESAT is slightly more general in that the ring $S$, the one  where the witness $\bff$ is typed over, is more general than in UCS. 

We choose to focus on UPESAT, rather than UCS, due to its greater simplicity and naturality. Further, on the one hand it is more general than UCS, and on the other, most of the territory in which UCS  is more general can be ported to UPESAT and our schemes in relatively straightforward manner, though technically involved. The exception being the support of $\QQ$ and $\QQ[X]$, whose inclusion in our work may require more than straightforward but technical work\footnote{This is because we commit to the bits that form our witnesses, and reconstruct the original witness from the bits via a linear map. At the very least, it is unclear how to bypass this gap if witnesses are rational \albert{explain somewhere that this is why we ask that rings are finitely generated as modules?}}.


%\medskip
We note that, in UCS, constraints are expressed through membership to ideals generated by monic non-constant polynomials. This is equivalent to expressing them directly over the quotient rings $R_i$ in UPESAT. 
%
%In this work, we will build a IOP for UPESAT with witnesses being committed over binary fields. The same ideas can be used to construct an analogous IOP for UCS with witnesses with entries in a bounded subset of $\ZZ[X]$ (as opposed to $\QQ[X]$ \albert{committing to rational through their bit-decomposition is a bit strange, but it may still be possible?}). 
 %
 %We choose to focus on UPESAT to favor simplicity of exposition and of our IOP. 
 %
%The main difference between UPESAT and UCS  is that the later allows for witnesses and constraints over arbitrary and multiple domains in the same instance, including rational polynomial rings $\QQ[X]$, and it allows to reuse witness elements in $\QQ[X]$ or $\ZZ[X]$ in other domains such as finite fields and polynomial rings over finite fields, by using   homomorphic projections.
\end{remark}






\subsection{An IOP for UPESAT with oracles in binary fields}\label{s:iop_upesat}
\albertimportant{Needs a pass to update notation}
%\john{Since we would have a lot of repeated hash computations, would there be a way to fold them like using sumfold before Step 2? Albert, I remember running an experiment on this using Hyrax and sumfold and seems to perform 2x faster than the baseline zinc+. Even though we used elliptic curve commitments, we maybe able to a similar folding technique } \albert{I thought about that. The issue is that the commitments are in binary fields and the claims to fold in $R[X]$, so folding here requires really novel research, if possible at all (efficiently).}

We now describe an IOP for $\UPESAT$. 
% 
%Given $a\in \ring[X]$ and $i\in [\iota]$, the equation $a=0$ holds in $R_i=\ring[X]/(h_i)$ if and only if $a$ belongs to the ideal $(h_i)$ generated by $h_i$ in $\ring[X]$. Accordingly, our IOP treats the UPESAT constraints as the ideal membership conditions $\pp_j^{(i)}(\xx,\bff)\in (h_i)$ for $j\in [\numrows]$ and $i \in [\iota]$.
% 
We assume $\numrows=2^{n'}$ is a power of two, and identify $[n]$ with $\bits^{n'}$.   Let $\codedim\geq 1$ and let $L\in S^{\codedim\times \codedim_0}$ be a \emph{bitification matrix} as in \cref{r: non_uniform_bit_decomposition}. Given $\bff\in \moduleonebound^{\codedim_0}$ we define $\bitstext(\bff)\in \bits^\codedim\subseteq \ZZ^\codedim$ so that $$L \cdot \bitstext(\bff)_S = \bff$$ (see \cref{r: non_uniform_bit_decomposition}). 
Let $\cC$ be a linear code over $\comfield$ of dimension $\codedim/\comfieldexp$ and blocklength $\codelength$. Let $\codedim_1,\codedim_2$ be tensor decomposition parameters, to be used in \cref{c:LinModFieldBin}, so that $\codedim=\codedim_1\cdot \codedim_2$. Let $\cC$ be  a linear code  over $\comfield$ of dimension $\codedim/\comfieldexp$. 


%For each prime, recall that $\pi_q:\ZZ\to \FF_q$ denotes reduction modulo $q$; and for every finite field $\KK$ of characteristic $q$,  we convene that $\pi_q$ acts on $\KK$ as the identity. 
%
%For each field $\KK$ of characteristic $q$, let $$\psi_\KK:\ZZ^{<\moduledegree}[X]\to \KK$$ be a $\pi_q$-module homomorphism \albert{give examples?}. If $\FF$ is another field of characteristic $q$, then we convene that $\psi_\KK$ is a $\pi_q=\mathsf{id}$-module homomorphism between the $\FF_q$-modules $\FF^{<\moduledegree}[X]$ and $\KK$. 
%
For each finite field $\KK$  let $\Pi_\KK$ be an IOR that given $(\inp,\bot, \wit)$ for  $\PESAT(\KK,\pp)$ (for any collection of polynomials $\pp$), outputs $(\inp',\bot, \wit) \in \LIN(\KK)$ (see \cref{?} for definitions of these relations). 
We emphasize that the output of $\Pi_\KK$ has the same witness $\wit$ as the input.
 
In our instantiations, $\Pi_{\KK}$ is  the usual Spartan \cite{spartan} PIOP over $\KK$.



\begin{construction}\label{c:iop_pimsat}
[IOP for {$\UPESAT$}]
\noindent \textbf{Parameters.} 
$(S,\codedim_0, B, \codedim, L,  \codedim_1,\codedim_2, \comfield,\cC)$ \albert{and $\bm{\psi}, \pp, \bm{R}$?} comprised of: a ring $S$; a witness length $\codedim_0$; a bit-size bound $B$; a bit amount $\codedim$; a bitification matrix $L$; a binary field $\comfield$; tensor decomposition dimensions $\codedim_1, \codedim_2$ with $\codedim_1\cdot \codedim_2= \codedim$; a linear code $\cC$ over $\comfield$ of dimension $\codedim/\comfieldexp$.\\


\noindent \textbf{Inputs.} The prover and verifier receive $\inp=\xx\in \moduleonebound^{\inplen}$. The prover additionally receives $\wit= \bff\in \moduleonebound^{\codedim_0}$ such that $(\inp,\bot, \wit) \in \UPESAT(\moduleonebound, \bm{\psi}, \pp, \bm{R})$. 


\begin{enumerate}[leftmargin=*] 
     \item \textbf{Round 0 (Optional): Out Of Domain (OOD) sample.} If the protocol is to be configured with proximity parameter $\delta$ beyond the unique decoding radius of $\cC$, the prover and verifier proceed exactly as in Round 0 of \cref{c:core_iop}. \albertimportant{Mention we can skip the OOD of \cref{c:core_iop} then}
\item \textbf{Round 1: Oracles.}
\begin{itemize}[leftmargin=*]
\item The prover  computes $$f= \Enc_{\cC}\circ \pack\circ\pi_2(\bitstext(\bff)) \in \comfield^{\codelength}$$  and sends  the verifier oracle access to $f$. 
\item The prover forms the multilinear polynomial, for all $i\in [\iota]$,
%
$$
Q^{(i)}(Z_1, \ldots, Z_\nu) \defeq \sum_{\bb\in\bits^{n'}}\pp_{\bb}^{(i)}(\psi_i(\xx),\psi_i(\bff)) \cdot \eq(\bb, Z_1, \ldots, Z_\nu) \in R_i[Z_1, \ldots, Z_\nu].
$$
%
The  claim that  $\mathbf{p}_{\bb}^{(i)}(\psi_i(\xx), \psi_i(\bff)) =0$ in $R_i$  for all  $\bb\in \bits^{n'}$ and $i \in [\iota]$ is equivalent to the claim 
%
\begin{equation}\label{e:iop_ideal_eq}
   Q^{(i)}(\zz) =0 \qquad  \text{ in } R_i \text{ for all } \zz\in \bits^{n'}, \ i \in [\iota].
\end{equation}
%
\item If $\ring=\FF$, we let $\KK$ be a large enough (to be determined later) extension of $\FF$. If $\ring=\ZZ$, verifier samples a prime  from a set of primes $\cP$ (cf.\ \cref{r:sampling_field} for comments on how to do this in practice). Abusing the notation, we denote the prime by $q$ and denote $\FF_q$ by $\KK$.

The claim \eqref{e:iop_ideal_eq} is replaced by
%
\begin{equation}\label{e:projected_Q_claim}
   \pi_q(Q^{(i)}(\zz)) \in (\pi_q(h_i)) \qquad  \text{ for all } \zz\in \bits^{n'}, \ i \in [\iota].
\end{equation}
Recall (cf.\ \cref{d:projections}) that $\pi_q$ is coefficient-wise reduction modulo $q$ when $\ring=\ZZ$, and it is the identity map when $\ring$ is a field $\FF$ of characteristic $q$.
\item The verifier samples and sends a random point $\rr\in \KK^{n'}$.
\end{itemize}
\item \textbf{Round 2: Ideal check}
\begin{itemize}[leftmargin=*]
\item The prover responds with values $u_i\in \KK[X]$ for each $i\in [\iota]$. Supposedly, 
%
$
u_i = \pi_q(Q^{(i)})(\rr).
$

These values can be computed as one would compute the corresponding values in a univariate skip, see \cref{s:univ_skip_ideal_membership} \albert{todo}. Alternatively, if $Q^{(i)}$ is linear on the multilinear extension of $\bff$, then  $u_i$ can easily be computed from $\mle{\pi_q(\bff)}(\rr)$, which in turn can be computed as $\moduledegree$ MLE evaluations over $\pi_q(\ring)$\footnote{This was already observed at least in \cite{flock}, when $Q^{(i)}$ is the multilinear extension of the vector $C\cdot \zz$ in an R1CS constraint \albert{check: is that linear on witness?}.}.
\item The verifier checks that $u_i \in (\pi_q(h_i))$ for all $i\in [\iota]$. Then it samples and sends batching elements $\theta_1, \ldots, \theta_\iota\in \KK$. Further, if $\moduledegree>0$, it samples and sends an element $\alpha$ \albert{it could send a random $\FF_q$-module homomorphism  $\KK^{<d}[X]\to \KK$}. 
\end{itemize}
\item \textbf{Using the IOR $\Pi_\KK$ for the projected claim}
\begin{itemize}[leftmargin=*]
\item Let $\pi\defeq \pi_{q,\alpha}$ (cf.\ \cref{d:projections}). The prover and verifier replace the claim \eqref{e:projected_Q_claim} with the claim 
%
\begin{equation}\label{e:batched_projected_claim}
\sum_{i\in [\iota]}\sum_{\bb\in \bits^{n'}}\theta_i\cdot \pi(\pp_\bb^{(i)})(\pi(\xx),\pi(\bff)) \cdot \eq(\bb, \rr) = \sum_{i\in [\iota]} \theta_i \cdot \pi(u_i) \in \KK.
\end{equation}
%
Observe that $(\pi(\xx),\bot, \pi(\bff))\in \PESAT(\KK,\hat{\pp})$ (cf.\ \cref{d:pesat})
with $\hat{\pp}=(\hat{\pp}_1)$ given by $$\hat{\pp}_1=\sum_{i\in[\iota]}\theta_i\cdot\sum_{\bb\in \bits^{n'}}\pi(\pp_{\bb}^{(i)})(Y_1, \ldots, Y_{\inplen+\codedim})-\sum_{i\in [\iota]} \theta_i\cdot \pi(u_i).$$ 
%
Accordingly, prover and verifier use the IOR $\Pi_{\KK}$  to reduce \eqref{e:batched_projected_claim} to 
 a claim of the form:
\begin{equation}\label{e:iop_linear_claim_raw}
((\vv,\mu), \bot, \pi(\bff))\in \LIN(\KK)
\end{equation}
%
for some $\vv \in \KK^{ \codedim_0}$  and some $\mu\in \KK$.
\end{itemize}
\item \textbf{Using the \ftwoz \ IOR for $\LinModFieldBin$.}
\begin{itemize}[leftmargin=*]
\item The initial claim $(\inp,\bot,\wit)\in \UPESAT$ has been replaced by $((\vv,\mu),\bot,\pi(\bff))\in \LIN(\KK)$.  Notice that $\pi:\ring[X] \to \KK$ induces a $\pi_q$-module homomorphism $\psi:\ringdegbound{\moduledegree}[X] \to \KK$ given by $\psi(h(X)) = \pi(h(X))$. Hence
%
\begin{equation}\label{e:iop_final_claim}
((\vv,\mu),f,\bff) \in \LinModFieldBin(\ringtwobounds{B}{\moduledegree},\psi,\KK, \cC).
\end{equation}
%
Prover and verifier now use \cref{c:LinModFieldBin} to prove  \eqref{e:iop_final_claim}, with parameters $$(R,q,\ringdegbound{\moduledegree}[X],\moduledegree,\KK, \psi, (1, X \ldots, X^{\moduledegree-1}), \codedim_0, B, \codedim_1,\codedim_2, \comfield,\cC_1^{\equiv \codedim_2}).$$
%
\end{itemize}
\end{enumerate}
    
\end{construction}


\begin{remark}[Randomly sampling a field]\label{r:sampling_field}
\albert{Add some comment about the recursive cost and the mitigations we have}
\end{remark}




\begin{theorem}[Completeness and round-by-round knowledge soundness of \cref{c:iop_pimsat}] \label{t:thm_iop_pimsat}  Let $(R,M_R, \moduledegree, \codedim_0, B, \codedim_1, \codedim_2, \comfield, \cC_1^{\equiv \codedim_2})$ be parameters for \cref{c:iop_pimsat}. Assume the following:
\begin{itemize}
\item For every field $\KK$ sampleable in Round 1 of \cref{c:iop_pimsat} (or for the fixed field $\KK$ if $\ring = \FF$), $\Pi_{\KK}$ is a $t_{\Pi_{\KK}}$-round perfectly complete IOR from $\PESAT(\KK, \pp)$ to $\LIN(\KK)$ with round-by-round knowledge soundness (without relaxation) and errors $\kappa_{\Pi_\KK, 1}, \ldots, \kappa_{\Pi_\KK, t_{\Pi_\KK}}$. 
\item Let $\delta\in (0,1)$ be a proximity parameter.
\item Assume $|\comfield|-1 \geq (\codedim_1+1)\cdot (q_{\max} -1)$, where $q_{\max}$ is the largest prime in $\cP$ in \cref{c:LinModFieldBin}.
\item Assume the notation and hypotheses of \cref{t:thm_core_IOPP,t:thm_LinModFieldBits}, so that  $\cref{c:LinModFieldBin}$ is a perfectly complete $(\tau_{\mathsf{ext}}+t_\red +t_\iopp)$-round IOP for $\LinModFieldBin(\ringtwobounds{B}{\moduledegree},\psi,\KK, \cC)$, with $\tau_{\mathsf{ext}}=1$ if the field $\KK$ in \cref{c:iop_pimsat} is a proper extension field and $\tau_{\mathsf{ext}}=0$ if $\KK$ is a prime field, with round-by-round knowledge with relaxation $(\LinModFieldBin(\ringtwobounds{B}{\moduledegree},\psi,\KK, \cC)^\delta,\bot)$ and errors $$\kappa_{\mathsf{lift-proj}}^*, (\ksred_i)_{i\in [t_\red]}, (\ksiopp_i)_{i\in [t_\iopp]},$$ where
    %
    $$
    \kappa_{\mathsf{lift-proj}}* = \begin{cases}\frac{\degevalfield}{q_{\min}} + \frac{\log(\codedim_1 +1) + \log(q-1)}{|\cP|\cdot \lfloor\log(q_{\min})\rfloor} & \text{if } \KK \text{ is a proper extension field,}\\
    \bot & \text{if } \KK \text{ is a prime field.}
    \end{cases}
    $$
    and $q_{\min}$ is the smallest prime in the set of sampleable primes $\cP$. 
\end{itemize}
    Then \cref{c:iop_pimsat} is a perfectly complete $(3+t_\KK + \tau_{\mathsf{ext}} + t_{\red} +t_\iopp)$-IOP for $\UPESAT(\ringtwobounds{B}{\moduledegree},\pp,\bm{R})$ with round-by-round knowledge soundness without relaxation and errors 
    %
    $$
    \kappa_{\mathsf{OOD}}, \kappa_1, \kappa_2, (\kappa_{\Pi_\KK,i})_{i\in [t_{\Pi_\KK}]}, \kappa_{\mathsf{lift-proj}}^*, (\ksred_i)_{i\in [t_\red]}, (\ksiopp_i)_{i\in [t_\iopp]}, 
    $$
    with 
    %
    \begin{align*}
    \kappa_1 &\defeq  \left(\frac{B_{\max}\cdot \left(1+\max\left(\nabla(\pp) - d_{\min},0\right)\right)}{|\cP|\cdot \lfloor\log(|\KK_{\min}|)\rfloor}\right)\cdot \tau_{\ZZ} + \frac{1}{|\KK_{\min}|} \\
    \kappa_2 &\defeq \frac{\nabla(\pp)}{|\KK_{\min}|} + \frac{1}{|\KK_{\min}|},
    \end{align*}
    %
    where:
    \begin{itemize}
        \item $\tau_{\ZZ}=1$ if $\ring=\ZZ$ and $\tau_\ZZ=0$ otherwise,
        \item $\KK_{\min}$ is either $\KK$ if $\ring=\FF$, or  the field $\FF_q$ with $q\in \cP$ with the smallest size if $\ring=\ZZ$,
        \item $d_{\min} \defeq \min\left\{\left.\deg(h_i) \ \right| \ i\in [\iota], \ h_i\neq 0\right\}$, with the convention $d_{\min}=+\infty$ when $h_i=0$ for all $i\in [\iota]$
        \item $\nabla(\pp)\defeq \max\left\{\max_{\bb\in \bits^{n'}, i\in[\iota]}\{\deg_X(\pp_{j}^{(i)}) + (\moduledegree-1)\cdot \max_{t\in [\inplen+\codedim\textit{}]}\deg_{Y_t}(\pp_j^i)\cdot (\inplen+\codedim)\}, \moduledegree\right\}$
        \item $B_{\max} \defeq \max_{\bb\in \bits^{n'}, i\in [\iota]} \lceil\log(\lVert\pp_{\bb}^{(i)}(\xx,\ff)\rVert_{\infty})\rceil. $
        \item         $\kappa_{\mathsf{OOD}}$ is as in \cref{t:thm_core_IOPP}.
    \end{itemize}
    %
    In \cref{t:errors_pimsat} we describe the variations in number of rounds and errors of \cref{c:iop_pimsat} depending on whether $\ring =\ZZ$ or $\ring=\FF$, and on whether $\KK$ is a prime field or a proper extension field.
\end{theorem}
\begin{table}[]\label{t:errors_pimsat}
    \centering
    \begin{tabular}{c|c|c|c}
        & $\ring=\ZZ$ & $\ring = \FF$, $\FF\subsetneq\KK$ &  $\ring = \FF=\KK$\newline \text{ is a prime field}\\
        \hline
        Rounds                   & $2+t_\KK + t_\red+t_\iopp$ & $2+t_\KK + 1+t_\red+t_\iopp$ & $2+t_\KK + t_\red+t_\iopp$ \\
        \hline
        $\kappa_{\mathsf{lift-proj}}$ & $\bot$ & $\frac{\degevalfield}{q_{\min}} + \frac{\log(\codedim_1 +1) + \log(q-1)}{|\cP|\cdot \lfloor\log(q_{\min})\rfloor}$ & $\bot$ \\
        \hline
        $\kappa_1$               & $\frac{B_{\max}\cdot \left(1+\max\left(\nabla(\pp) - d_{\min},0\right)\right)}{|\cP|\cdot \lfloor\log(|\KK_{\min}|)\rfloor} + \frac{1}{|\KK_{\min}|}$ & $\frac{1}{|\KK_{\min}|}$ & $\frac{1}{|\KK_{\min}|}$
    \end{tabular}
    \caption{\albertimportant{Outdated. Fix.} Variations in number of rounds and errors of \cref{c:iop_pimsat} depending on whether $\ring =\ZZ$ or $\ring=\FF$, and on whether $\KK$ is a prime field or a proper extension field}
\end{table}


We defer the proof of \cref{t:thm_iop_pimsat} to \cref{a:iop_pimsat} \albertimportant{Errors after composition of IOR's cannot be the same, since they receive different inputs?}. We state and prove a key lemma here:
\begin{lemma}[Ideal membership is preserved under modular reduction]\label{lem:ideals_modular projection}
Let $d,B\ge 1$ and $f,h\in\mathbb{Z}_{\leq B}^{<d}[X]$, with $h$ either $0$ or a nonzero monic polynomial. Let $\cP$ be a set of primes. Then
%
\[
\Pr_{p\leftarrow\mathcal{P}}
\Big[\;\pi_p(f)\in (\pi_p(h))\ \wedge\ f\notin (h)\;\Big]
\;<\;\frac{B\cdot (1 +\max(\deg(f)-\deg(h)+1,0))}{|\mathcal{P}|\cdot \lfloor\log(q_{\min})\rfloor}.
\]
where $q_{\min}\defeq \min \cP$.
\end{lemma}
\begin{proof}
Assume  $f\not\in(h)$. In particular
$f\neq 0$. 
 If $h=0$, then  $(\pi_p(h))=\{0\}$, and so
 $\pi_p(f)=0$. Since $\|f\|_\infty\leq 2^{B}-1$, there are at most $B/\lfloor \log(q_{\min})\rfloor$ primes $p$ such that $\pi_p(f)=0$, as needed.

If $h\neq0$ and $\deg(h)=0$, then $h=1$ (since $h$ is monic) and so $f\in(h)=\ZZ[X]$, a contradiction. Assume thus that $\deg(h)\geq 1$. Let $Q,R\in\mathbb{Z}[X]$ be such that
$f=h\cdot Q+R$ with $\deg(R)<\deg(h)$ (such \emph{integer} polynomials exist because $h$ is monic). Since
$f\notin(h)$, we have $R\neq0$.
By \cref{?}, $$\|R\|_\infty \le\|f\|_\infty\cdot (1 +\|h\|_\infty)^{\max(\deg(f) - \deg(h) +1,0)}.$$ 
Now since $\deg(\pi_p(h))=\deg(h)$ (because $h$ is monic) we have $ \deg(\pi_p(h))> \deg(R)\geq \deg(\pi_p(R))$.  The condition $\pi_p(f)\in (\pi_p(h))$ thus implies that $\pi_p(R)=0$.  There are at most $\log_2(\|R\|_\infty)/\lfloor \log q_{\min} \rfloor$ primes dividing the coefficients of $R$. Hence
%
$$\Pr_{p\gets \cP}[f\not\in (h) \ \wedge \pi_p(f)\in (\pi_p(h))] \leq \frac{\log_2(\|R\|_\infty)}{|\cP|\cdot \lfloor \log q_{\min} \rfloor} \leq \frac{B + \max(\deg(f)-\deg(h)+1,0)\cdot B}{|\cP|\cdot \lfloor \log q_{\min}\rfloor},$$
as needed.
%
\end{proof}




\albert{Ideal membership as a generalization of univariate skip?}\label{s:univ_skip_ideal_membership}
\albert{Repeated computations through ideal membership?}



\subsection{Completeness and soundness of \cref{c:iop_pimsat}}\label{a:iop_pimsat}



To prove \cref{t:thm_iop_pimsat}, we will need to introduce an extension of the definition of $\PESAT(\KK, \pp)$ and $\LIN(\KK)$ so that these relations can hold witnesses in $\ringtwobounds{B}{\moduledegree}$ and their encodings over a binary field. Namely, given a field of characteristic $q$ and a $\pi_q$-module homomorphism (cf.\ \cref{?}) $\psi_\KK:\ringdegbound{\moduledegree}[X]\to \KK$, we define:
$$
\PESAT(\ringtwobounds{B}{\moduledegree}, \KK, \pp,\psi_\KK, \cC) \defeq \left\{ \left(\begin{aligned} &\inp=\xx\in \KK^\inplen, \\&\oracle = f \in (\comfield)^{\codelength}, \\ &\wit= \bff\in (\ringtwobounds{B}{\moduledegree})^{\codedim_0}\end{aligned}\right) \;\middle|\;  \begin{aligned} 
&f=\Enc_{\cC}\circ \pack \circ \pi_2(\bitstext(\bff)) \\ \wedge \  &\pp_i(\xx, \psi_\KK(\bff)) = 0  \text{ for all } i\in [\numrows]\end{aligned}\right\}
$$
Note that when $\pp$ contains a single linear polynomial, then
%
$$
\PESAT(\ringtwobounds{B}{\moduledegree}, \KK, \pp, \modulemorphism, \cC) = \LinModFieldBin(\ringtwobounds{B}{\moduledegree},\modulemorphism, \KK, \cC).
$$
%
%We define $\LIN$ similarly: $\LIN(\ringtwobounds{B}{\moduledegree}, \KK, \psi_\KK, \cC) $ is simply the above $\PESAT$ with $\pp$ containing a single linear polynomial on $\xx, \psi_\KK(\bff)$.



We will also need an encoded version of UPESAT:
we define    $\UPESAT(\ringtwobounds{B}{\moduledegree},\pp,\bm{R}, \cC)$ as 
$$
\UPESAT(\ringtwobounds{B}{\moduledegree}, \pp,\bm{R}, \cC) 
%\UPESAT(\ringtwobounds{B}{\moduledegree}, \pp,\bm{R}) 
= \left\{\left. 
\begin{aligned}
    \inp &= \xx \in \left(\ringtwobounds{B}{\moduledegree}[X]\right)^{\inplen},\\
    \oracle &= f \in \comfield^{\codelength},\\
    \wit &=\bff \in \left(\ringtwobounds{B}{\moduledegree}[X]\right)^\codedim
\end{aligned} 
\ \right| \ 
\begin{aligned}
    & f = \Enc_{\cC}\circ \pack \circ \picomfield(\bitstext(\bff)),\\
    &\mathbf{p}_{j}^{(i)}(\xx, \bff) = 0 \ \text{ in } R_i \ \text{ for all } j\in [\numrows], \ i \in [\iota]
\end{aligned}
\right\}
$$
%For ease of notation we write 
%
%$$
%\UPESAT(\ringtwobounds{B}{\moduledegree}, \pp,\bm{R})  = \UPESAT, \qquad \UPESAT(\ringtwobounds{B}{\moduledegree}, \pp,\bm{R}, \cC)  = \UPESAT(\ringtwobounds{B}{\moduledegree}, \pp,\bm{R}, \cC).
%$$
%
As usual, for the relations $\Relation$ that contain the constraint $f=\Enc_{\cC}\circ \pack \circ \pi_2(\bitstext(\bff))$, we also consider their \emph{relaxed} version $\Relation^\delta$, where such constraint is replaced by $\hamming(f,\Enc_{\cC}\circ \pack \circ \pi_2(\bitstext(\bff)))<\delta$.






\begin{lemma}\label{l:lift_pesat_to_encodings}
Let $\KK$ be a finite field of characteristic $q$ and let $\psi_\KK:\ringdegbound{\moduledegree}[X]\to \KK$ be a $\pi_q$-module homomorphism \albert{correct terminology? And above?}. Let $\Pi_\KK$ be a $t_\KK$-round perfectly complete IOR from $\PESAT(\KK,\pp)$ to $\LIN(\KK)$ with round-by-round knowledge soundness (with no relaxation) with errors $(\kappa_{\Pi_\KK, i})_{i\in [t_\KK]}$. 

Let $\delta\in (0,1)$. Then there is an IOR $\Pi_\KK'$ from $\PESAT(\ringtwobounds{B}{\moduledegree}, \KK,\pp,\psi_\KK, \cC)$ to $\LinModFieldBin(\ringtwobounds{B}{\moduledegree},\psi_\KK, \KK,\cC)$  with $t_\KK$ rounds, perfectly complete, and with round-by-round knowledge soundness with relaxation $(\PESAT(\ringtwobounds{B}{\moduledegree}, \KK,\pp,\psi_\KK, \cC)^\delta,\LinModFieldBin(\ringtwobounds{B}{\moduledegree},\psi_\KK), \KK, \cC)^\delta)$ and with errors $(\kappa_{\Pi_\KK, i})_{i\in [t_\KK]}$.  

\albertimportant{Could there be an issue with extracting the correct bit size $B$? No because prover outputs $\bff$, not $\bff'$} This IOR $\Pi_\KK'$ is exactly $\Pi_\KK$, except that prover and verifier receive an oracle $f$, which they do not use throughout the IOR; and additionally the prover receives a witness $\bff$ with entries in $\ringtwobounds{B}{\moduledegree}$ instead of in $\KK$. Then the prover of $\Pi_\KK'$ proceeds exactly as the prover of $\Pi_\KK$ after replacing $\bff$ by $\psi_\KK(\bff)$. At the end of the IOR, both prover and verifier of $\Pi_\KK'$ output the same  $\inp$ output by $\Pi_\KK$. Additionally, they output $\oracle=f$, and the prover of $\Pi_\KK'$ outputs any vector in the preimage of $\psi_\KK^{-1}(\bff')$, where $\bff'$ is the witness output by the prover of $\Pi_\KK$. 
\end{lemma}
\begin{proof}
\albertimportant{TODO. Anywhere in the literature?}
\end{proof}

\bigskip

\noindent We build  three IORs $\Pi_1,\Pi_2,\Pi_3$ whose composition closely mimics \cref{c:iop_pimsat}.  We discuss the case when Round 0 does not trigger (i.e.\ $\delta$ is within the unique decoding raidus). The case when it does follows in the same way as in \cref{thm:?} (see the proof in \cref{?}), and uses known arguments from \cite{stir,whir}.

%The first one, $\Pi_0$ \albertimportant{Not needed?}, is a $1$-round IOR from 


%\paragraph{Description of $\Pi_1$.} In $\Pi_0$, given $(\inp,\bot, \wit)\in \PIMSAT(\ringtwobounds{B}{\moduledegree},\pp,\bm{\ideal})$ as input, the prover sends $f=\Enc_{\cC}\circ \pack \pi_2(\bitstext(\bff))$ to the verifier, and the verifier replies with an empty message. Then the IOR outputs $(\inp,f,\wit)\in \PIMSAT(\ringtwobounds{B}{\moduledegree},\pp,\bm{\ideal}, \cC, \pack)$. The first item in Round 1 of \cref{c:iop_pimsat} follows precisely this IOR: prover computes and sends $f$, and prover and verifier are left with $(\inp,f)$ and $(\inp,f,\bff)\in \PIMSAT(\cC, \pack)$.




\paragraph{Description of $\Pi_1$.} The IOR $\Pi_1$ is a two-round IOR from $\UPESAT(\ringtwobounds{B}{\moduledegree}, \pp,\bm{R}, \cC)$ to $\PESAT(\ringtwobounds{B}{\moduledegree},\KK, \hat{\pp},\pi,\cC)$, for some polynomials $\hat{\pp}$ and $\pi_q$-module homomorphism $\pi:\ringtwobounds{B}{\moduledegree} \to \KK$. In $\Pi_1$, verifier and prover receive $(\inp, f)$ and $(\inp, f,\bff)\in \UPESAT(\ringtwobounds{B}{\moduledegree}, \pp,\bm{R}, \cC)$ as inputs, respectively. In the first round, prover and verifier simulate Round 1 of \cref{c:iop_pimsat}, skipping the first item (computing and sending the encoding $f$ of $\bff$), since $f$ is already in the input of the prover and verifier. The prover instead  sends an empty message $\bot$. For the second round of $\Pi_1$, prover and verifier simulate Round 2 of \cref{c:iop_pimsat}. At the end of $\Pi_1$, prover and verifier output $$(\pi(\xx),f,\bff)\in \PESAT(\ringtwobounds{B}{\moduledegree},\KK, \hat{\pp},\pi,\cC),$$
where $\hat{\pp}$ and $\pi$ are determined during the execution of $\Pi_1$ \albert{This may be formally incorrect: I don't think current difinitions support relations in an IOR changing depending on the run of the IOR, so we'd need to add the parameters in $\inp$}.
%

\paragraph{Description of $\Pi_2$.} 
The IOR $\Pi_2$ is an IOR from $\PESAT(\ringtwobounds{B}{\moduledegree},\KK, \hat{\pp},\pi,\cC)$ to $\LIN(\ringtwobounds{B}{\moduledegree},\KK, \hat{\pp},\cC)$. Prover and verifier receive $(\pi(\xx),f)$ and $(\pi(\xx),f,\bff)$ as inputs. Then they  apply the IOR
$\Pi_\KK'$ given by \cref{l:lift_pesat_to_encodings}, where $\Pi_\KK$ is the IOR applied in Item 3 of \cref{c:iop_pimsat}. They use as inputs $(\pi(\xx),f)$ and $(\pi(\xx),f,\bff)$, respectively. By \cref{l:lift_pesat_to_encodings}, $\Pi_\KK'$ is an IOR from $\PESAT(\ringtwobounds{B}{\moduledegree},\KK, \hat{\pp},\pi, \cC)$ to $\LIN(\ringtwobounds{B}{\moduledegree},\KK,\pi,\cC)$, perfectly complete, with $t_\KK$ rounds, and with round-by-round knowledge with relaxation  $(\PESAT(\ringtwobounds{B}{\moduledegree},\KK, \hat{\pp}, \pi,\cC)^\delta,\LIN(\ringtwobounds{B}{\moduledegree},\KK,\pi,\cC)^\delta)$ and with errors $(\kappa_{\Pi_\KK,i})_{i\in t_{\KK}}$ 

\paragraph{Description of $\Pi_3$}

\albertimportant{Should it be named $\Pi_4$?}
The IOR $\Pi_3$ is an IOP for $\LIN(\ringtwobounds{B}{\moduledegree},\KK,\pi,\cC)$. It consists  in applying Phase 4 of \cref{c:iop_pimsat}, i.e.\ simulating \cref{c:LinModFieldBin} with inputs $((\vv,\mu), f,\bff)$. By \cref{t:thm_LinModFieldBits} and the hypothesis of \cref{t:thm_iop_pimsat}, \cref{c:LinModFieldBin} is a $(\tau_{\mathsf{\mathsf{ext}}}+t_\red+t_\iopp)$-round perfectly complete IOP for $\LIN(\ringtwobounds{B}{\moduledegree},\KK,\pi,\cC)$ with round-by-round knowledge with relaxation $(\LIN(\ringtwobounds{B}{\moduledegree},\KK,\pi,\cC)^\delta,\bot)$, with errors $\kappa_{\mathsf{lift-proj}}^*, (\ksred_i)_{i\in [t_\red]}, (\ksiopp_i)_{i\in [t_\iopp]}$.


The next lemma states that the composition of $\Pi_1, \Pi_2$ and $\Pi_3$ is an IOR whose completeness and soundness implies completeness and soundness of \cref{c:iop_pimsat}.

\begin{lemma}\label{l:Pi_123_is_good}
Let $\Pi_{123}$ be the composition of $\Pi_1,\Pi_2$ and $\Pi_3$.  Then $\Pi_{123}$ is a $(2+t_\KK + \tau_{\mathsf{ext}} +t_\red + t_\iopp)$-round IOP for $\UPESAT(\ringtwobounds{B}{\moduledegree}, \pp,\bm{R}, \cC)$. Assume $\Pi_{123}$ is perfectly complete and, for any $\delta\in (0,1)$,  has round-by-round knowledge with relaxation $(\UPESAT(\ringtwobounds{B}{\moduledegree}, \pp,\bm{R}, \cC)^\delta, \bot)$ with errors $(\kappa_j')_{j\in [2+t_\KK + \tau_{\mathsf{ext}} +t_\red + t_\iopp]}$. 

Then \cref{c:iop_pimsat} is a $(2+t_\KK + \tau_{\mathsf{ext}} +t_\red + t_\iopp)$-round perfectly complete IOP for $\UPESAT(\ringtwobounds{B}{\moduledegree}, \pp,\bm{R})$,  and has  round-by-round knowledge without relaxation with errors $(\kappa_j)_{j\in [2+t_\KK + \tau_{\mathsf{ext}} +t_\red + t_\iopp]}$ defined as
%
$$
\kappa_j(\inp) \defeq \max_{f \text{ s.t. } (\inp,f)\not\in \UPESAT(\ringtwobounds{B}{\moduledegree}, \pp,\bm{R}, \cC)^\delta} \kappa_j(\inp, f)'.
$$
\albertimportant{Need to argue somewhere that these errors are negligible if $\kappa_j'$ are -- I'd prefer to just use the errors from $\kappa_j'$ directly. This may be related with the issue about composition of IORs and having ill-defined errors} 
\end{lemma}
\begin{proof}
    Deferred to \cref{a:Pi_123_is_good}.
\end{proof}

Due to \cref{l:Pi_123_is_good}, it  suffices to show that $\Pi_{123}$ is perfectly complete and has the desired soundness properties.  We prove this by using the composition of IOR result \cref{r:ior_compositoin}.  Both $\Pi_2$ and $\Pi_3$ already have the desired properties by definition and hypotheses of \cref{t:thm_core_IOPP}. Thus it suffices to show the following:
%
\begin{lemma}[Completeness and soundness of $\Pi_1$]\label{l:completeness_soundness_Pi_1}
    Let $\delta\in (0,1)$. The IOR $\Pi_1$ is a perfectly complete $2$-round IOR from $\UPESAT(\ringtwobounds{B}{\moduledegree}, \pp,\bm{R}, \cC)$ to $\PESAT(\ringtwobounds{B}{\moduledegree}, \KK,\hat{\pp}, \pi,  \cC)$ with round-by-round knowledge soundness  with relaxation $(\UPESAT(\ringtwobounds{B}{\moduledegree}, \pp,\bm{R}, \cC)^\delta,\PESAT(\ringtwobounds{B}{\moduledegree}, \KK,\hat{\pp}, \pi,  \cC)^\delta)$ and with errors 
    \begin{align*}
    \kappa_1 &=  \frac{B_{\max}\cdot \left(1+\max\left(\nabla(\pp) - d_{\min},0\right)\right)}{|\cP|\cdot \lfloor\log(|\KK_{\min}|)\rfloor} + \frac{1}{|\KK_{\min}|} \\
    \kappa_2 &= \frac{\nabla(\pp)}{|\KK_{\min}|} + \frac{1}{|\KK_{\min}|},
    \end{align*}
\end{lemma}
\begin{proof}
\albertimportant{Probably needs some polishing and review}
    Indeed, let $\inp=\xx, \oracle=f$ be an input and oracle for $\UPESAT(\ringtwobounds{B}{\moduledegree}, \pp,\bm{R}, \cC)$. For any $\ww$ define $\kstate(\inp,f, \ww)=1$ if and only if $(\inp,f,\ww)\in \UPESAT(\ringtwobounds{B}{\moduledegree}, \pp,\bm{R}, \cC)^\delta$. Given a full transcript $\tr=(\bot, (\KK,\rr),  (u_1,\ldots, u_\iota),(\theta_1,\ldots, \theta_\iota, \alpha))$ for $\Pi_{1}$, define $\kstate(\inp, f, \tr, \ww) = 1$ if and only if $(\pi(\xx), f, \pi(\bff))\in \PESAT(\ringtwobounds{B}{\moduledegree}, \KK, \hat{\pp},  \pi, \cC)^\delta$, where $\pi=\pi_{q,\alpha}$, $q$ being the characteristic of $\KK$ and $\alpha$ one of the verifier's messages in $\tr$. For partial transcripts where the prover $\tr$ is about to send a message $m$, let $\kstate(\inp,f,\tr, m, \ww) = \kstate(\inp, f,\tr,\ww)$. 

    We now define  $\kstate(\inp, f, \tr, \ww)$ where $\tr$ is nonempty and such that prover is about to move.  The only undefined setting is when $\tr=(\bot, (\KK,\rr))$. In this case we define $\kstate(\inp, f, \tr, \ww) = 1$ if and only if $$\Delta(f,\Enc_\cC\circ \pack \circ \pi_2(\bitstext(\ww)) < \delta$$ and $\pi_q(Q_{\ww}^{(i)})(\rr)\in (\pi_q(h_i))$ for all $i\in [\iota]$,  where 
    %
    $$
    Q_{\ww}^{(i)}(\mathbf{Z})\defeq\sum_{\bb\in\bits^{n'}}\pp_{\bb}^{(i)}(\xx,\ww) \cdot \eq(\bb, \mathbf{Z}) \in \ring[X]. 
$$
%
Observe that $\kstate$ satisfies the three conditions of \cref{def:knowledge-state-function}.

  %Let $\tr$ be a transcript on this input as  $\cE_{\mathsf{rbr}}(\inp, f,\bot,(\KK,\rr), (u_i)_{i\in [\iota]},\ww)=\ww$.
      
    %If $\tr=(\bot)$ and verifier is about to send $(\KK,\rr)$, we define $\kstate(\inp, f, \bot, \ww)=\kstate(\inp, f,\ww)$. We define the extractor on this input as  $\cE_{\mathsf{rbr}}(\inp, f,\bot,\ww)=\ww$.   

 We define the extractor  $\cE_{\mathsf{rbr}}$ as $\cE_{\mathsf{rbr}}(\inp,f,\tr,\ww)=\ww$ for all $\inp,f,\tr,\ww$. Let $\tr$ be a transcript for $\Pi_1$ where the verifier is about to move, so that either $\tr=(\bot)$ or $\tr=(\bot, (\KK,\rr), (u_i)_{i\in[\iota]})$. For a transcript $\tr=(\bot)$, we need to bound the following probability, corresponding to the round-by-round knowledge soundness of Round 1 of $\Pi_1$:
%
\begin{align*}
  P_1\defeq &\Pr_{\KK,\rr}\!\left[
  \exists\,\kswit \ \left| \ \begin{array}{l}
  \wedge\ \kstate\bigl(\inp,f,\bot,\ww\bigr)=0\\
  \wedge\ \kstate(\inp,f, \bot,(\KK,\rr),\ww)=1
  \end{array}
  \right.\right] 
\end{align*}
% 
By our definition of $\kstate$, $P_1$ is the probability that 
%
\begin{equation}\label{e:boring_proof_eq1}\pi_q(Q_{\ww}^{(i)})(\rr) \in (\pi_q(h_i)) \text{ for all } i\in [\iota] \text{ and } \Delta(f, \Enc_{\cC}\circ \pack \circ \pi_2(\bitstext(\ww))<\delta
\end{equation} while at the same time \begin{equation}\label{e:boring_proof_eq2}\pp_\bb^{(i)}(\xx,\ww)\not\in (h_i) \text{ for some } \bb\in \bits^{n'} \text{ and } i\in [\iota]\end{equation} or $\Delta(f, \Enc_{\cC}\circ \pack \circ \pi_2(\bitstext(\ww))\geq\delta$. Clearly, the latter condition cannot hold if \eqref{e:boring_proof_eq1} holds; hence we will study the probability  that \eqref{e:boring_proof_eq1} and \eqref{e:boring_proof_eq2} hold simultaneously.  We ultimately want to prove that $P_1 \leq \kappa_1$.

Let $E_1$ be the event that \eqref{e:boring_proof_eq2} holds and $\pi_q(\pp_\bb^{(i)}(\xx,\ww))\in (\pi_q(h_i))$ for all $\bb,i$. By \cref{lem:ideals_modular projection}, 
%
$$
\Pr_{\KK,\rr}[E_1] \leq \max_{\bb\in \bits^{n'},i\in[\iota]}\left\{\frac{B'\cdot (1 +\max(\deg(\pp_{\bb}^{i}(\xx,\bff))-\deg(h_{i})+1,0))}{|\mathcal{P}|\cdot \lfloor\log(q_{\min})\rfloor}\right\},
$$
where $\deg(h_{i})\defeq+\infty$ if $h_i=0$; and $B'$ is $\lceil\log(\|\pp_{\bb}^{(i)} (\xx,\bff)\|_\infty)\rceil$. We have, with $d_{\min}$ as in \cref{t:thm_iop_pimsat}, \albert{more detail?}
%
\begin{align*}
&\max_{\bb\in \bits^{n'},i\in[\iota]}\left\{\deg(\pp_{\bb}^{i}(\xx,\bff))-\deg(h_i)\right\} \\  &\qquad\qquad \qquad \leq\nabla(\pp) -d_{\min}= \max_{\bb\in \bits^{n'}, i\in[\iota]}\{\deg_X(\pp_{\bb}^i) + (\moduledegree-1)\cdot \max_{t\in [\inplen+\codedim\textit{}]}\deg_{Y_t}(\pp_{\bb}^i)\} + \moduledegree -d_{\min},\\
& \max_{\bb\in \bits^{n'}, i\in [\iota]} B' = \max_{\bb\in \bits^{n'}, i\in [\iota]} \lceil\log(\lVert\pp_{\bb}^{(i)}(\xx,\ff)\rVert_{\infty})\rceil = B_{\max}.
\end{align*}
%
Now, let $E$ be the event that both \eqref{e:boring_proof_eq1} and \eqref{e:boring_proof_eq2} hold. Then
%
$$
P_1=\Pr_{\KK,\rr}[E] \leq \Pr_{\KK, \rr}[E \ | \ \neg E_1] + \Pr_{\KK,\rr}[E_1].
$$
%
We next analyze $\Pr_{\KK, \rr}[E \ | \ \neg E_1]$. If $E_1$ does not occur, then the  multilinear polynomial $\pi_q(Q_\ww^{(i)})(\ZZZ)$ is not the zero polynomial in $\KK[X]$. On the other hand, for $E$ to occur, we need  $\pi_q(Q_\ww^{(i)})(\rr)=0$ in $\KK$. Hence, by Schwartz-Zippel lemma,  $\Pr_{\KK, \rr}[E \ | \ \neg E_1] \leq 1/|\KK_{\min}|$. It follows that 
%
$
P_1 \leq \kappa_1,
$ as needed.

Finally, we bound the following probability, corresponding to the round-by-round knowledge soundness of the second round of $\Pi_1$.
\begin{align*}
  P_2\defeq &\Pr_{(\theta_i)_{i\in[\iota], \alpha}}\!\left[
  \exists\,\kswit \ \left| \ \begin{array}{l}
  \wedge\ \kstate\bigl(\inp,f,\bot,(\KK,\rr), (u_i)_{i\in[\iota]},\ww\bigr)=0\\
  \wedge\ \kstate(\inp,f, \bot,(\KK,\rr),(u_i)_{i\in [\iota]},((\theta_i)_{i\in[\iota]},\alpha), \ww)=1
  \end{array}
  \right.\right] 
\end{align*}
%
By definition, \begin{equation}\label{e:boring_proof_eq3}\kstate(\inp,f, \bot,(\KK,\rr),(u_i)_{i\in [\iota]},((\theta_i)_{i\in[\iota]},\alpha), \ww)=1 \ \Leftrightarrow \ (\pi(\inp), f, \pi(f))\in \PESAT(\ringtwobounds{B}{\moduledegree}, \KK, \hat{\pp}, \cC)^\delta\end{equation} and verifier does not abort \albert{It's not completely transparent how verifier accepts or aborts in the definition of RBR KS}; while $$\kstate\bigl(\inp,f,\bot,(\KK,\rr), (u_i)_{i\in[\iota]},\ww\bigr)=\kstate\bigl(\inp,f,\bot,(\KK,\rr),\ww\bigr)=0$$  if and only if 
$\Delta(f,\Enc_\cC\circ \pack \circ \pi_2(\bitstext(\ww)) \geq \delta$ or \begin{equation}\label{e:boring_proof_eq4}\pi_q(Q_{\ww}^{(i)})(\rr)\not\in (\pi_q(h_i)) \text{ for some } i\in [\iota].\end{equation}  
%
Let $E$ be the event that  \eqref{e:boring_proof_eq3} and \eqref{e:boring_proof_eq4} hold. Again, $P_2= \Pr_{(\theta_i),\alpha}[E]$, since \eqref{e:boring_proof_eq3} forces $\Delta(f,\Enc_\cC\circ \pack \circ \pi_2(\bitstext(\ww)) < \delta$.

Now, if $E$ occurs then
$$
\sum_{i\in [\iota]}\sum_{\bb\in \bits^{n'}}\theta_i\cdot \pi(\pp_\bb^{(i)})(\pi(\xx),\pi(\ww)) \cdot \eq(\bb, \rr) = \sum_{i\in [\iota]} \theta_i \cdot \pi(u_i) \in \KK,
$$
and at the same time 
$$
\pi_q(Q_{\ww}^{(i)})(\rr) = \sum_{\bb\in\bits^{n'}}\pi_q(\pp_{\bb}^{(i)}(\xx,\bff)) \cdot \eq(\bb,\rr) \not\in (\pi_q(h_i))\subseteq \KK[X]
$$
 for some $i\in [\iota]$. Since the verifier does not reject, $u_i \in (\pi_q(h_i))$ for all $i\in [\iota]$. Hence, if $E$ holds, then for some $i$, we have $\pi_q(Q_{\ww}^{(i)}(\rr)) \neq u_i$. In particular $u_i' \defeq \pi_q(Q_{\ww}^{(i)}(\rr)) - u_i \in \KK[X]$ is a nonzero polynomial with coefficients in $\KK$.

Let $E_1$ be the event that  $E$ holds and
%
$\pi(Q_{\ww}^{(i)}(\rr)) =\pi(u_i)$ for all $i\in [\iota]$ (recall that $\pi=\pi_{q,\alpha}$ acts on $\KK[X]$ by replacing $X$ by $\alpha$).  Now, since $u_i'(X)$ is a nonzero polynomial from $\KK[X]$, we have that $\pi(u_i'(X)) = u_i'(\alpha) =0$ with probability at most $\deg(u_i')/|\KK_{\min}|$ by Schwartz-Zippel lemma. It follows that \albertimportant{Again, I'm not missing a union bound right?}
%
%
\begin{align*}
\Pr_{(\theta_i)_i, \alpha}[E_1] &\leq \max_{i\in [\iota]} \frac{\deg(\pi_q(Q_{\ww}^{(i)}(\rr)) - u_i)}{|\KK_{\min}|}\leq  \max_{i\in [\iota]} \frac{\max\left\{\deg(\pi_q(Q_{\ww}^{(i)}(\rr))),\moduledegree\right\}}{|\KK_{\min}|}\\ &\leq  \max_{\bb\in \bits^{n'},i\in [\iota]} \frac{\max\left\{\deg(\pp_\bb^{(i)}(\xx,\ff)),\moduledegree\right\}}{|\KK_{\min}|} \leq  \frac{\nabla(\pp)}{|\KK_{\min}|},
\end{align*}
%
where we used that $\deg_X(u_i)< \moduledegree$; that $\pi_q(Q_{\ww}^{(i)}(\rr))\in \KK[X]$ is a $\KK$-linear combination of the polynomials $\pi_q(\pp_{\bb}^{(i)}(\xx,\bff)) \in \KK[X]$ \albert{More details about the last inequality?}. 
%
Now,
$$
\Pr[E]\leq \Pr[E\mid \neg E_1] + \Pr[E_1].
$$
%
If $E_1$ does not hold, but $E$ holds, then $\pi(Q_{\ww}^{(i)}(\rr)) \neq \pi(u_i)$ for some $i\in [\iota]$. Then the polynomial  on variables $U_1,\ldots, U_\iota$,

%
$$
Q'(U_1,\ldots, U_\iota)\defeq \sum_{i\in[\iota]} U_i \cdot \pi(Q_{\ww}^{(i)}) - \sum_{i\in [\iota]}U_i \cdot \pi(u_i)
$$
%
is a nonzero multilinear polynomial with coefficients in $\KK$. Now, $[E\mid \neg E_1]$ holds if and only if $Q'(\theta_1, \ldots, \theta_\iota)=0\in \KK$. By Schwartz-Zippel lemma, this probability is at most $1/|\KK_{\min}|$.  It follows that $P_2\leq \kappa_2$, as required. 
%
\end{proof}
This completes the proof of \cref{t:thm_iop_pimsat}.

%Rounds 1 and 2 of \cref{c:iop_pimsat} follow this IOR: the only difference between the rounds of \cref{c:iop_pimsat} and $\Pi_1$ is the specification of the input and output of the IOR. Besides those, the prover and verifier operate exactly the same. Further, in \cref{c:iop_pimsat}, the prover and verifier have both the data that comprises the input and output to $\Pi_1$. \albert{Is this enough or should I make the $\Pi_i$ prover and verifier simulate provers and verifiers of \cref{c:iop_pimsat}?}



%---
%and verifier  defined as Rounds 1 and 2 of \cref{c:iop_pimsat} after the first item (where prover sends $f$), and formally we have the prover send an empty message in $\Pi_2$ as its first message.  $\Pi_2$ takes as input $(\inp, f, \ww)\in \UPESAT(\ringtwobounds{B}{\moduledegree}, \pp,\bm{R}, \cC)$, and outputs $(\pi(\xx),f, \pi(\bff))\in \PESAT(\KK, \hat{\pp}, \cC)$. We will argue later that $\Pi_2$ is a perfectly complete $2$-round IOR from $\UPESAT(\ringtwobounds{B}{\moduledegree}, \pp,\bm{R}, \cC)$ to $\PESAT(\ringtwobounds{B}{\moduledegree},\KK, \hat{\pp}, \pi, \cC)$ \albert{Remark that $\pi$ is a module homomorphism}, with round-by-round knowledge soundness with relaxation $(\UPESAT(\ringtwobounds{B}{\moduledegree}, \pp,\bm{R}, \cC)^\delta,\PESAT(\ringtwobounds{B}{\moduledegree},\KK, \hat{\pp}, \pi, \cC)^\delta)$ and errors $\kappa_1, \kappa_2$ as in the statement of the present theorem.  \albert{Need to argue better that Rounds 1 and 2 can be seen as executing this IOR.}

%Once we have established our claims regarding $\Pi_2$, the result then follows from  the composition of IORs result (cf.\ \cref{prop:composed}).
%


\subsubsection{Proof of \cref{l:Pi_123_is_good}}\label{a:Pi_123_is_good}

 We prove \cref{l:Pi_123_is_good}. Observe that the prover and verifier of $\Pi_{123}$ proceed in exactly the same way as the prover and verifier of \cref{c:iop_pimsat} after the prover sends $f$ in Item 1 of Round 1 (except that in $\Pi_{123}$ the prover sends an initial empty message).  

If $\Pi_{123}$ is perfectly complete, then clearly \cref{c:iop_pimsat} is as well \albert{more details?}. Further,  transcripts between $\Pi_{123}$ and \cref{c:iop_pimsat} are the same, except that the prover in $\Pi_{123}$ sends an empty first message when the prover in \cref{c:iop_pimsat} sends the encoding $f$. More precisely, a (partial) nonempty transcript of \cref{c:core_iop} has the form $(f,\tr')$ for some $\tr'$ and some first prover message $f$, and the ``mirror'' transcript of $\Pi_{123}$ is $(\bot,\tr')$, with $f$ being now part of the instance oracle.

From any knowledge state function with relaxation $\kstate_{\Pi_{123}}$ for $\Pi_{123}$, we  define a state function $\kstate$ for \cref{c:iop_pimsat} as follows: let $\tr$ be a partial  or full transcript of \cref{c:iop_pimsat}. If $\tr$ is not empty, write $\tr=(f,\tr')$ as above. Then we define
%
$$
\kstate(\inp,\bot, (f,\tr'), \ww) = \kstate_{\Pi_{123}}(\inp, f, (\bot, \tr'), \ww).$$ 
%
If $\tr$ is empty, we define
%
$$
\kstate(\inp,\bot, \ww) = 1 \ \Leftrightarrow \ (\inp,\bot,\ww) \in \UPESAT.
$$
%
Hence $\kstate$ satisfies Condition 1 of \cref{def:knowledge-state-function} (Empty transcript). We now show it satisfies Conditions 2 (Prover moves) and 3 (Full transcript), assuming $\kstate_{\Pi_{123}}$ satisfies \cref{def:knowledge-state-function}. 

Let $\tr$ be a transcript of \cref{c:iop_pimsat} where the prover is about to move and $\kstate(\inp,\bot,\tr,\ww)=0$. If $\tr$ is nonempty, then $\tr=(f,\tr')$ as above, and  for any next prover message $m$ of \cref{c:iop_pimsat}, $\kstate(\inp,\bot, (\tr, m), \ww)= \kstate_{\Pi_{123}}(\inp,f, (\bot, \tr', m),\ww)=0$ because, since $\kstate_{\Pi_{123}}$ satisfies \cref{def:knowledge-state-function}, $\kstate_{\Pi_{123}}(\inp,f, (\bot, \tr'),\ww)=\kstate(\inp,\bot, (f, \tr'),\ww)=0$. If $\tr$ is empty, then the next prover's message is $f$. If $\kstate(\inp,\bot,\ww)=0$, then $(\inp,\bot,\ww)\not\in \UPESAT$. Now $\kstate(\inp,\bot, f, \ww) = \kstate_{\Pi_{123}}(\inp, f, \bot,\ww)$. By the "Prover moves" condition of \cref{def:knowledge-state-function}, if $\kstate_{\Pi_{123}}(\inp, f,\ww)=0$ (empty transcript where the prover is about to move), then $\kstate_{\Pi_{123}}(\inp, f, \bot,\ww)=0$ ($\bot$ is always the first message of the prover in $\Pi_{123}$). Now $\kstate_{\Pi_{123}}(\inp, f,\ww)=0$ indeed because by the condition ``Empty transcript'' of \cref{def:knowledge-state-function}, $\kstate_{\Pi_{123}}(\inp, f,\ww)=0$ if and only if $(\inp, f,\ww)\not\in \UPESAT(\ringtwobounds{B}{\moduledegree}, \pp,\bm{R}, \cC)^\delta$, and since $(\inp,\bot,\ww)\not\in \UPESAT$, it follows that $(\inp, f,\ww)\not\in \UPESAT(\ringtwobounds{B}{\moduledegree}, \pp,\bm{R}, \cC)^\delta$.  Hence $\kstate(\inp,\bot, f, \ww)=0$ and $\kstate$ satisfies the condition ``Prover moves'' of \cref{c:iop_pimsat}.
 

 Regarding the third condition, ``Full transcript'', let $(f,\tr')$ be a full transcript for \cref{c:iop_pimsat}.  Then for all $\inp,\ww$,
%
$$
\kstate(\inp,\bot,(f,\tr'),\ww) = \kstate_{\Pi_{123}}(\inp,f,(\bot,\tr'), \ww) = 1 \ \Leftrightarrow \ \verifier_{\Pi_{123}}(\inp,f, (\bot,\tr'))=1, 
$$
%
where $\verifier_{\Pi_{123}}$ is the verifier of $\Pi_{123}$.

Notice that the verifier of \cref{c:iop_pimsat} given $(\inp, \bot, (f,\tr'))$ makes exactly the same checks as $\verifier_{\Pi_{123}}(\inp,f, (\bot,\tr'))$. Hence the verifier of \cref{c:iop_pimsat} accepts $(\inp, \bot, (f,\tr'))$ if and only if $\kstate(\inp,\bot,(f,\tr'),\ww) = 1$, as needed. This shows that $\kstate$ satisfies the three conditions of \cref{def:knowledge-state-function}.




Now assume $\Pi_{123}$ has round-by-round knowledge soundness with relaxation $(\UPESAT(\ringtwobounds{B}{\moduledegree}, \pp,\bm{R}, \cC)^\delta,\bot)$ with errors $(\kappa_{j}')_{j\in [2+t_\KK, 1+ t_\red +t_\iopp]}$, with respect to the state function $\kstate_{\Pi_{123}}$. We claim that then \cref{c:iop_pimsat} has round-by-round knowledge soundness without relaxation, with errors $(\kappa_{j}')_{j\in [2+t_\KK, 1+ t_\red +t_\iopp]}$, with respect to the state function $\kstate$.

We start by defining an extractor $\cE$ for \cref{c:iop_pimsat} simply as $$\cE(\inp,\bot, (f,\tr'),\ww)= \cE_{\Pi_{123}}(\inp,f, (\bot,\tr'), \ww)$$ for every (partial) transcript $(f,\tr')$ where the verifier is about to move. We need not define an extractor for, say $(\inp, \bot, \ww)$, since the verifier is not about to move after an empty transcript $\bot$ (the prover is about to send $f$).
%

%
Now, let $\tr=(f,\tr')$ be a (partial)  transcript of \cref{c:iop_pimsat} where the verifier is about to send its $j$-th message $\rho$. Then we have, using the definitions of $\kstate, \cE$ and our assumption on the round-by-round knowledge soundness of $\Pi_{123}$,
%
\begin{align*}
  &\Pr_{\verifiermsg}\!\left[
  \exists\,\kswit \ \left| \ \begin{array}{l}
  \wedge\ \kstate\bigl(\inp,\bot,(f,\tr'),\cE(\inp,\bot, (f,\tr'),\verifiermsg,\ww)\bigr)=0\\
  \wedge\ \kstate(\inp,\bot, (f,\tr'),\verifiermsg,\ww)=1
  \end{array}
  \right.\right]\\ = 
  &\Pr_{\verifiermsg}\!\left[
  \exists\,\kswit \ \left| \ \begin{array}{l}
  \wedge\ \kstate_{\Pi_{123}}\bigl(\inp,f,(\bot,\tr'),\cE_{\Pi_{123}}(\inp,f, (\bot,\tr'),\verifiermsg,\ww)\bigr)=0\\
  \wedge\ \kstate_{\Pi_{123}}(\inp,f,(\bot,\tr'),\verifiermsg,\ww)=1
  \end{array}
  \right.\right] \\ \le&\ \knowledgeerror_j(\inp, f)' \leq \max_{f \text{ s.t. } (\inp,f)\not\in \UPESAT(\ringtwobounds{B}{\moduledegree}, \pp,\bm{R}, \cC)^\delta} = \kappa_j(\inp).
\end{align*}
%
This concludes the proof of \cref{l:Pi_123_is_good}.
%(\inp, f, \ww)\in \UPESAT(\ringtwobounds{B}{\moduledegree}, \pp,\bm{R}, \cC)^\delta \ \Rightarrow \ (\inp, \bot, \ww) \in \UPESAT. 



